\section{求解流程与结果分析（待填充）}

\subsection{求解算法框架}
建议采用\u201c构造启发式 + 局部搜索\u201d两阶段流程：
\begin{enumerate}[label=\arabic*)]
    \item 初始解构造：聚类/扫描法 + 最近邻插入；
    \item 邻域改进：2-opt、跨路径交换、重定位；
    \item 约束修复：时间窗与容量违规惩罚回退；
    \item 输出结果：路径、成本、排放、服务质量指标。
\end{enumerate}

\subsection{结果表格模板}
\begin{table}[H]
\centering
\caption{三类问题结果汇总（示例模板）}
\begin{tabular}{cccccc}
\toprule
场景 & 总成本(元) & 碳排放(kg) & 车辆数(辆) & 平均迟到(h) & 求解时间(s) \\
\midrule
问题一 & 待填 & 待填 & 待填 & 待填 & 待填 \\
问题二 & 待填 & 待填 & 待填 & 待填 & 待填 \\
问题三 & 待填 & 待填 & 待填 & 待填 & 待填 \\
\bottomrule
\end{tabular}
\end{table}

\subsection{图件插入模板}
\begin{figure}[H]
    \centering
    \includegraphics[width=0.8\textwidth]{route_example.png}
    \caption{路径规划结果示意图（占位图，待替换）}
\end{figure}

\subsection{敏感性分析建议}
建议至少分析以下敏感性：
\begin{enumerate}[label=\arabic*)]
    \item 碳排放成本系数变化对车型选择的影响；
    \item 迟到惩罚系数变化对时间窗满足率的影响；
    \item 绿色区半径或限行时段变化对总成本的影响。
\end{enumerate}
